\section{Relation to the literature}\label{sec:11}
We separate what is used from what is proved.

\textbf{Classical inputs, used and not claimed.} Two are external: the decidability of integer quadratic solvability \cite{GS04}, and DPRM \cite{DPR61,Mat70}. \cref{sec:10} uses a third, the eventual representation theorem for ternary lattice cosets of Kane and Kim \cite{KK26}, through \cref{lemma:10.8}. Skolem's degree reduction, and with it $PP_{2} = SE_{4} = \mathrm{Dioph}$, is proved here rather than assumed (\cref{theorem:6.2}); $\mathrm{REC} ⊊ \mathrm{RE}$ and the equality $\mathrm{TotComp} = \mathrm{Fun}(PP_{2})$ under DPRM are derived in Section 6. \cref{sec:6} states each of these as a named hypothesis or proves it from mathlib's computability library. Dirichlet's theorem on primes in arithmetic progressions and Lagrange's four-square theorem are used from mathlib.

\textbf{Contributions.} (1) The uniform finite-fibre quadratic height bound (\cref{theorem:2.1}); (2) the finite-fibre cardinality bound with the same constant (\cref{theorem:2.2}); (3) quadratic growth of SE₂-definable total functions (\cref{theorem:3.2}); (4) the sharp injective dimension theorem $n ≤ m$, its converse, and the structural alternative that yields it (\cref{theorem:4.1}, \cref{theorem:4.2}, \cref{theorem:4.3}), with $\mathrm{ca}(SE_{2}) = 1$ as a corollary; (5) the classification $\mathrm{graph}(p) ∈ SE_{2} ⟺ \mathrm{deg}\ p ≤ 2$ (\cref{theorem:5.1}); (6) non-closure of Fun(SE₂) under composition, and exactness of $\mathrm{Fun}(SE_{2}) ⊊ \mathrm{TotComp}$ by arity (\cref{theorem:3.4}, \cref{theorem:3.5}); (7) the typed certificate language with per-kind soundness (\cref{definition:8.2}, \cref{theorem:8.3}); (8) the unconditional structured summit (\cref{theorem:8.5}); (9) the arithmetic projection dichotomy in both its bounded and its data-carrying form, for every number of unknowns (\cref{theorem:9.2}, \cref{theorem:9.3}); (10) the effective diagonalization lemma for coded families (\cref{theorem:6.5}) and its consequence at arity one (\cref{corollary:6.6}); (11) the classification of SE₂-definable total functions under $\mathrm{TCER}$ (\cref{proposition:10.9,proposition:10.14}), with converse constructions and the case without unknowns that need no analytic input (\cref{proposition:10.16}).

The correct summary of the paper is this: the structural theorems identify the internal finite-fibre and functional constraints on the single-quadratic fragment, and combining them with classical decidability, DPRM and degree reduction places the fragment exactly in the computability hierarchy. The phase-diagram arrows themselves are not claimed as new; what is new is the pair of constraints that make the diagram exact, and the positive description of the infinite fibres.

\textbf{Comparison with adjacent programmes.} The existential-rank literature (Daans–Dittmann–Fehm \cite{DDF} and related work) bounds the \emph{number} of existential variables while leaving the degree free, for the domains and the classes of definable sets considered there; the present fragment bounds the \emph{degree} while leaving the number of unknowns free. The degree–unknowns trade-offs for universal equations (Jones \cite{Jon82}: every r.e. set by one equation of fixed degree in a fixed number of unknowns, degree and unknowns traded against each other) are the classical form of that tension; the finite-fold programme (Matiyasevich \cite{Mat10}) asks that each parameter point have finitely many \emph{witnesses}, a bound on the witness side, whereas our finite fibres are on the \emph{output} side with the witnesses unrestricted. We prove no refinement in either direction between these hierarchies and ours, and the objects studied differ: fibre growth, function growth and graph classification do not appear there. \cref{theorem:5.1} is not derived from the degree-free arity collapse of [DDF, Lemmas 8.19–8.20], and that collapse is not derived here from \cref{theorem:5.1}; whether either implies the other is a question this paper does not settle. The closest published statement to our \cref{sec:8} is in \cite{GS04}. Besides deciding solvability, the algorithm there decides, for a single given equation and a region cut out by finitely many rational affine hyperplanes, whether the solution set is finite, lists it when it is, and reports infinitude when it is not. Two differences should be stated plainly. First, that is an algorithmic statement about one instance: it answers a question about the equation in front of it, and it does not by itself yield a constant valid at every parameter point, which is what \cref{theorem:2.1} asserts and what the whole of Sections 3–5 rests on. Second, it reports infinitude rather than describing it, whereas \cref{theorem:8.5} exhibits the infinite fibre as generated by data of one of three named kinds. We do not claim that our results could not be derived from theirs with further work; we claim that the statements are different, and that the uniformity in the parameters and the structural description of the infinite case are what this paper adds.

\textbf{Presburger-definable functions.} Presburger cell decomposition describes the definable functions on $ℤ^{n}$ by finitely many linear formulas on cells cut out by linear inequalities and congruences (Cluckers \cite{Clu03}, Definitions 1–2 and Theorem 1). \cref{sec:10} concerns a different restriction: the graph of a total function must be the projection of the zero set of one polynomial of total degree at most two. With two or more outputs the classified functions are Presburger-definable and form a proper subclass: their linear parts coincide on all residue classes, and $a ↦ (|a|, 0)$ is Presburger-definable without a common slope. With one output the common-slope and parity-affine families are Presburger-definable, while a function with a genuinely quadratic part, such as $a ↦ a^{2}$, is not; so the scalar class is not contained in the Presburger-definable functions. The contrast with Presburger arithmetic does not come from an inability of one quadratic equation to express inequalities: four squares define the relation $a ≥ 0$. It comes from asking for the graph of a total function, with a unique output at every input. The one-hot quadratic penalty behind the converse constructions is a standard device of quadratic unconstrained optimization; the characterization of the functions is what \cref{sec:10} adds.

\textbf{The nearest neighbour on the infinite side.} Grechuk and Agbanwa \cite{GA26} fix a non-degenerate quadratic form F and a univariate integer polynomial P and develop a method for showing that P(x) is represented by F for infinitely many integers x. The method covers certain cubic and quartic P, and certain P of the form R(Q(x)) with $\mathrm{deg} R = 3$ and $\mathrm{deg} Q = 2$; the headline application is that x⁶ − 4 is a sum of two squares infinitely often, hence that $y^{2} + x^{3}y + z^{2} + 1 = 0$ has infinitely many integer solutions. Their main results are formalized in Lean. This is the nearest neighbour to \cref{theorem:8.5} in the current literature and deserves to be placed precisely. It is a 2026 arXiv preprint, and we cite it as such: nothing below treats it as a published theorem. Their question concerns one pair (F, P) at a time and is answered by establishing infinitude; ours fixes a single equation of total degree at most two and asserts, with one constant valid at every parameter point, that each scalar fibre is either bounded or carries arithmetic data generating an infinite output subset. Neither statement contains the other: they obtain infinitude for families we do not treat, and we obtain uniformity in the parameter together with a description of the infinite case, which they do not claim. Their own reading of the boundary agrees with ours: they observe that when $\mathrm{deg} P ≤ 2$ the resulting equation is in principle covered by the Grunewald–Segal algorithm, which is the same placement of \cite{GS04} we make in \cref{definition:1.5} and in the comparison above. We do not use their method and make no priority claim against it. That both developments are Lean-formalized is a coincidence of practice, not a mathematical relation. Set against each other at the level of statements:

\begin{longtable}{@{}>{\raggedright\arraybackslash}p{0.187\linewidth}>{\raggedright\arraybackslash}p{0.338\linewidth}>{\raggedright\arraybackslash}p{0.375\linewidth}@{}}
\toprule
 & Grechuk–Agbanwa \cite{GA26} & this paper \\
\midrule\endhead
object & one fixed non-degenerate quadratic form F and one univariate P, of certain degrees 3 and 4 and certain composites R∘Q & one equation of total degree ≤ 2, with parameters \\
question & is P(x) represented by F for infinitely many integers x? & for every parameter point, is the scalar fibre bounded, and if not, by what data is it generated? \\
quantifier over parameters & none: the pair (F, P) is fixed & one constant C chosen after the equation and before the parameter \\
conclusion in the infinite case & infinitude is established & an arithmetic certificate of one of three named kinds is exhibited \\
logical application & not the aim & finite-fibre relations, definable functions, the computability placement of SE₂ \\
Lean scope & the main results & the whole structure theory, Sections 2–9 \\
status & ar\allowbreak{}Xiv preprint, July 2026 & — \\
\bottomrule
\end{longtable}
Neither row of conclusions implies the other: they treat representation problems we do not, and we obtain parameter-uniformity together with a description of the infinite case, which they do not claim.

The theory of integral binary quadratic forms and their unit orbits (Lagrange reduction, Nagell \cite{Nag}) is an \emph{input} to \cref{sec:9} and to \cref{sec:10} rather than a competitor: what is done here with it is the uniformity in the affine data, which is what the projection dichotomy adds. The quartic collapse $SE_{4} = \mathrm{Dioph}$ is classical and is quoted as such.

\textbf{Search bounds.} The nearest quantitative neighbours of \cref{theorem:2.1} are the search bounds for Diophantine equations, a line beginning with Cassels \cite{Cas55} — an isotropic integral quadratic form in n variables has a nonzero integral zero of size at most a power of its largest coefficient — and continued by many authors, including Masser \cite{Mas98} for the rational case; Fukshansky \cite{Fuk13} surveys the direction. The closest single statement is Dietmann \cite{Die03}, as printed in §6 of the survey \cite{Fuk13} (with the $s = 2$ case due to Kornhauser): for an inhomogeneous integral quadratic Q in $s ≥ 3$ variables with nonsingular quadratic part and coefficient height at most H, there is a polynomial bound $Λ_{s}$(H) such that if Q has an integer zero at all, it has one of size at most $Λ_{s}$(H). Dietmann's theorem, as stated in Fukshansky's peer-reviewed survey, gives a small zero under existence assumptions and a coefficient-height bound; the original Dietmann article was not independently audited here, and it is not a load-bearing input of this paper. Three differences should be stated. First, that literature bounds the SMALLEST solution, whereas \cref{theorem:2.1} bounds EVERY element of a fibre that is finite; neither quantifier implies the other, since a finite fibre may have a small least element and a large greatest one. Second, their bound is a function of the coefficient height of the one equation at hand, whereas ours fixes the equation with a parameter block and asserts a single constant valid at every parameter point, with an exponent that does not grow with the number of unknowns. Third, Dietmann assumes $s ≥ 3$ and a nonsingular quadratic part, whereas \cref{theorem:2.1} assumes only total degree at most two — the singular and low-arity cases, which are exactly where \cref{sec:9} does its work, are inside our statement. We do not claim independence from this line, and it is a fair question whether a uniform-in-parameters bound can be extracted from those methods; we claim that the statements differ in the three ways above.

